\documentclass[11pt]{article}
\usepackage[a4paper,margin=1in]{geometry}
\usepackage{fontspec}
\usepackage[boldfont=false]{xeCJK}
\setCJKmainfont{FandolSong-Regular.otf}
\usepackage{amsmath}
\usepackage{amssymb}
\usepackage{array}
\usepackage{booktabs}
\usepackage{graphicx}
\usepackage{adjustbox}
\usepackage{enumitem}
\setlist[itemize]{topsep=2pt,partopsep=0pt,parsep=0pt,itemsep=2pt,leftmargin=1.4em}
\setlist[enumerate]{topsep=2pt,partopsep=0pt,parsep=0pt,itemsep=2pt,leftmargin=1.6em}
\usepackage{parskip}
\usepackage{fvextra}
\fvset{breaklines=true,breakanywhere=true,fontsize=\small}
\setlength{\parindent}{0pt}

\begin{document}

\begin{center}
  {\LARGE\bfseries Chemical energetics}\\[0.35em]
  {\large Igcse Chemistry}
\end{center}
\vspace{0.6em}

\section*{Exothermic and endothermic reactions}

In every chemical reaction, energy is transferred. The reaction is either exothermic or endothermic, depending on which way the energy moves.

\begin{itemize}
\item An \textbf{exothermic reaction} 放热反应 gives out \textbf{thermal energy} 热能 to the \textbf{surroundings} 环境. So the temperature of the surroundings goes \textbf{up}.

\item An \textbf{endothermic reaction} 吸热反应 takes in thermal energy from the surroundings. So the temperature of the surroundings goes \textbf{down}.

\end{itemize}

Examples of exothermic reactions are \textbf{combustion} 燃烧 (burning a fuel) and \textbf{neutralisation} 中和 (an acid reacting with an alkali). An example of an endothermic reaction is the thermal \textbf{decomposition} 分解 of a compound (breaking it down using heat).

\section*{Enthalpy change, ΔH}

The amount of thermal energy transferred in a reaction is called the \textbf{enthalpy change} 焓变. It is written as $\Delta H$ and has these signs:

\begin{itemize}
\item $\Delta H$ is \textbf{negative} for an exothermic reaction, because energy leaves the chemicals.

\item $\Delta H$ is \textbf{positive} for an endothermic reaction, because energy is taken in.

\end{itemize}

\section*{Activation energy}

Particles do not react every time they meet. The \textbf{activation energy} 活化能 ($E_a$) is the smallest amount of energy that \textbf{colliding} 碰撞 \textbf{particles} 粒子 must have before they can react. It is like a hill the particles must get over before the reaction can happen.

\section*{Reaction pathway diagrams}

A \textbf{reaction pathway diagram} 反应进程图 shows how the energy changes as a reaction happens. Energy is on the vertical axis, and the progress of the reaction is on the horizontal axis.

\begin{itemize}
\item The line starts at the energy level of the \textbf{reactants} 反应物.

\item It rises over a 'hill' — the height of this hill is the activation energy $E_a$.

\item It then falls or rises to the energy level of the \textbf{products} 生成物.

\item The gap between the reactant level and the product level is the enthalpy change $\Delta H$.

\end{itemize}

For an \textbf{exothermic} reaction, the products are \textbf{lower} than the reactants, so energy is given out and $\Delta H$ is negative.

For an \textbf{endothermic} reaction, the products are \textbf{higher} than the reactants, so energy is taken in and $\Delta H$ is positive.

\section*{Bonds and energy}

A reaction involves breaking the bonds in the reactants and making new bonds in the products.

\begin{itemize}
\item \textbf{Bond breaking} 断键 takes in energy, so it is an endothermic step.

\item \textbf{Bond making} 成键 gives out energy, so it is an exothermic step.

\end{itemize}

The \textbf{bond energy} 键能 is the energy needed to break one mole of a particular bond. You can use bond energies to find the enthalpy change:

\[
\Delta H = (\text{energy to break all bonds}) - (\text{energy released making all bonds})
\]

If more energy is given out making bonds than is taken in breaking bonds, the reaction is exothermic ($\Delta H$ negative). If less energy is given out, it is endothermic ($\Delta H$ positive).

\subsection*{Worked example}

For the reaction $\text{H}_2 + \text{Cl}_2 \rightarrow 2\text{HCl}$, use these bond energies (in kJ/mol): H–H $= 436$, Cl–Cl $= 242$, H–Cl $= 431$.

Bonds broken: one H–H and one Cl–Cl $= 436 + 242 = 678$.

Bonds made: two H–Cl $= 2 \times 431 = 862$.

\[
\Delta H = 678 - 862 = -184 \text{ kJ/mol}
\]

The answer is negative, so this reaction is exothermic.


\end{document}
